\chapter{ClaimSimulator}
\label{chap:claim_simulator}

\section{Purpose}

The \texttt{ClaimSimulator} class generates synthetic claim-level data
from a policy table. While \texttt{StochasticSimulator} produces aggregate
losses for a single distributional specification, \texttt{ClaimSimulator}
is designed for portfolios of policies that may have different exposure
windows and different frequency or severity assumptions.

This is the right tool when the analyst needs:

\begin{itemize}[leftmargin=*]
    \item A claim-level dataset with policy identifiers and occurrence
          dates.
    \item Synthetic data suitable for testing reserving methods.
    \item A starting point for downstream development analysis (e.g.\
          loss triangles).
\end{itemize}

\section{Required Policy Schema}

The class accepts a \texttt{pandas.DataFrame} with one row per policy.
The required columns are listed below.

\begin{center}
\begin{tabular}{ll}
\toprule
\textbf{Column} & \textbf{Description} \\
\midrule
\texttt{policy\_id}   & Unique identifier for the policy. \\
\texttt{freq\_dist}   & Name of the frequency distribution. \\
\texttt{freq\_params} & Tuple of frequency distribution parameters. \\
\texttt{sev\_dist}    & Name of the severity distribution. \\
\texttt{sev\_params}  & Tuple of severity distribution parameters. \\
\texttt{start\_date}  & Policy effective date. \\
\texttt{end\_date}    & Policy expiration date. \\
\bottomrule
\end{tabular}
\end{center}

The constructor validates that all required columns are present and
that \texttt{freq\_params} and \texttt{sev\_params} are tuples.

\section{Construction}

A typical policy table and \texttt{ClaimSimulator} construction looks
like this:

\begin{lstlisting}[style=pythonstyle]
import pandas as pd
import numpy as np
from actsim import ClaimSimulator

policies = pd.DataFrame({
    "policy_id":  range(1, 101),
    "freq_dist":  "poisson",
    "freq_params": list(zip(np.random.uniform(0.6, 0.8, 100).round(2),)),
    "sev_dist":   "lognormal",
    "sev_params": list(zip(
        np.random.uniform(8, 12, 100).round(2),
        np.random.uniform(0.3, 0.7, 100).round(2),
    )),
    "start_date": pd.Timestamp("2023-01-01"),
    "end_date":   pd.Timestamp("2023-12-31"),
})

claim_sim = ClaimSimulator(policies, random_seed=42)
\end{lstlisting}

Optional arguments include \texttt{correlation}, \texttt{copula\_type},
and \texttt{copula\_param} for introducing dependency between frequency
and severity within each policy group.

\section{Simulating Claims}

The \texttt{simulate\_claims} method groups policies by their distribution
specifications, runs a separate \texttt{StochasticSimulator} for each
group, and assembles the results into a unified claim DataFrame.

\begin{lstlisting}[style=pythonstyle]
claim_sim.simulate_claims()
print(claim_sim.claim_data.head())
\end{lstlisting}

The resulting \texttt{claim\_data} DataFrame contains the columns
\texttt{year}, \texttt{event\_id}, \texttt{yearly\_event\_id},
\texttt{amount}, \texttt{policy\_id}, \texttt{start\_date}, and
\texttt{end\_date}. The \texttt{year} column corresponds to the policy
index within its group rather than a calendar year, and is used
internally to map simulated claims back to their parent policy. This 
design reflects the package's original focus on catastrophe modelling, 
and the column naming and structure will be refined in a future version.

\subsection{Policies with Zero Simulated Claims}

It is possible for a policy to produce zero simulated claims, especially
when the expected frequency is below one. In that case, the policy
remains in the \texttt{policies} table but does not appear in
\texttt{claim\_data}. \actsim does not, by default, create explicit
zero-claim records. If downstream analyses require a record per policy
regardless of claim activity, the user should join \texttt{claim\_data}
back onto \texttt{policies} and treat missing rows as zero.

\section{Simulating Claim Dates}

After claims have been simulated, occurrence dates can be assigned using
a non-homogeneous Poisson process via the \texttt{simulate\_dates\_nhpp}
method.

\begin{lstlisting}[style=pythonstyle]
claim_sim.simulate_dates_nhpp(
    lambda0=10,
    alpha=0.5,
    phase=0,
    T=1,
)
\end{lstlisting}

The arguments control the seasonal intensity:
\begin{itemize}[leftmargin=*]
    \item \texttt{lambda0} -- baseline intensity.
    \item \texttt{alpha} -- seasonality amplitude.
    \item \texttt{phase} -- phase shift of the seasonality.
    \item \texttt{T} -- duration of the exposure period in years.
\end{itemize}

For each claim, the method draws a fraction of the exposure period from
the NHPP and maps that fraction to a date within the relevant policy's
active window. The amount column is renamed to \texttt{ultimate\_loss}
in preparation for the development step.

\section{Simulating Claim Development}

The \texttt{simulate\_claim\_development} method develops each claim
backwards from its ultimate amount, using a user-supplied set of base
loss development factors and a stochastic volatility multiplier.

\begin{lstlisting}[style=pythonstyle]
base_LDFs = {
    0:  2.00,
    3:  1.50,
    6:  1.20,
    9:  1.10,
    12: 1.05,
    15: 1.02,
    18: 1.00,
}

claim_sim.simulate_claim_development(
    base_LDFs=base_LDFs,
    volatility=0.1,
    cumulative_factor=1.0,
)
\end{lstlisting}

For each claim, a unique set of stochastic LDFs is drawn by multiplying
each base factor by a normal random variable with mean $1$ and standard
deviation \texttt{volatility}. Cumulative development factors (CDFs) are
constructed by chaining the stochastic LDFs from the latest age backwards.
The incurred loss at age $a$ is then computed as
\[
    L_a = \frac{U}{\mathrm{CDF}_a},
\]
where $U$ is the ultimate loss for the claim.

\subsection{Output Schema}

The \texttt{claim\_development} attribute is a long-format DataFrame
with one row per (claim, development month). The columns are listed
below.

\begin{center}
\begin{tabular}{ll}
\toprule
\textbf{Column} & \textbf{Description} \\
\midrule
\texttt{accident\_year}     & Accident year of the claim's incurred date. \\
\texttt{incurred\_date}     & Date the claim occurred. \\
\texttt{claim\_id}          & Unique claim identifier. \\
\texttt{policy\_id}         & Identifier of the parent policy. \\
\texttt{development\_month} & Development age in months. \\
\texttt{incurred\_loss}     & Incurred loss at this age. \\
\texttt{development\_date}  & Calendar date corresponding to the
                              development age. \\
\bottomrule
\end{tabular}
\end{center}

\section{Saving Claim Development Data}

Simulated development data can be persisted to CSV using the
\texttt{save\_claim\_development} method:

\begin{lstlisting}[style=pythonstyle]
claim_sim.save_claim_development("output/claim_development.csv")
\end{lstlisting}

The resulting file is suitable for ingestion by triangle-based reserving
tools, including the \texttt{chainladder} package.
